% page 146 du fac-simile (image p-145 du scan)
% bloc 160.0 x 237.7 mm ; marge gauche 8.0 mm
\pgc{-12.700mm}{0.000mm}{
\l{\cel{4}136.}
\l{}
\l{}
\l{\sou{enorma}.\cel{2}Qua ecesas omna grandeso o quanteso mezurebla. - DEF}
\l{}
\l{\sou{enoyar}.\cel{2}(\sou{netrans}.)\cel{2}Esar la objekto di ta deskomforto quan}
\l{\cel{3}sentas la psiko kande nulo interesas, okupas lu. - EFI.}
\l{}
\l{\sou{enrejistrar}. (\sou{trans.}) I. Enskribar o transkribar oficale en}
\l{\cel{3}registri. - II. Dicesas anke pri ula aparati qui notas}
\l{\cel{3}ocili, movi diversa (kom ex. la sfigmografo enrejistras}
\l{\cel{3}la pulso arteriala). - F.}
\l{}
\l{\sou{ensemblo}.\cel{2}Asemblajo de personi, de kozi, qui omna formacas}
\l{\cel{3}toto. - DEFR.}
\l{}
\l{\sou{ento}. - To quo existas. - EFI.}
\l{}
\l{\sou{entablamento}. (\sou{arkitekt.})\cel{2}Saliajo an la suprajo di la muregi}
\l{\cel{3}di edifico, sur qua su apogas la karpenturo di la tekto.-}
\l{\cel{3}DEFIRS.}
\l{}
\l{\sou{entamar}\cel{2}(\sou{trans.)}\cel{2}I. (\sou{pri kozo tillore netushita)}. Agar la}
\l{\cel{3}trancho, frapo, brecho, nocho ye lu, dat-unesme. - II.}
\l{\cel{3}(\sou{metaf.}) Komencar (ulo), iniciar, inaugurar. - F.}
\l{}
\l{\sou{enterito}.\cel{2}(\sou{patol.}) Influro di la mukozo qua tapete kovras}
\l{\cel{3}la kanalo intestina. - EF.}
\l{}
\l{\sou{enternar}.\cel{2}(\sou{trans}.)Koaktar a rezido en ta o ca urbo, regiono,}
\l{\cel{3}la ekiro esante interdiktita. - DEFIS.}
\l{}
\l{\sou{entimemo}. (\sou{logiko}). Silogismo en qua un ek la premisi esas}
\l{\cel{3}tacita, pro ke lua evidenteso posibligas suplear lu. -}
\l{\cel{3}DEFIRS.}
\l{}
\l{\sou{entomofilo}. (Planto) di qua la fekundigeso devenas de la trans-}
\l{\cel{3}porteso, da la insekti, di la poleno.}
\l{}
\l{\sou{entomologio}. Parto di la zoologio, qua traktas la insekti.-}
\l{\cel{3}DEFIRS.}
\l{}
\l{\sou{entraprezar}. (\sou{trans.}) I. Asumar la direkto, la exekuto (di}
\l{\cel{3}afero, di tasko). - II. Obligar su exekutor ula labori,}
\l{\cel{3}furnisor ulo, po preco o segun kondicioni determinita.--FIS.}
\l{}
\l{\sou{entratenar}. (\sou{trans.}) Igar provizata ye la kozi qui esas necesa}
\l{\cel{3}por vivo e por komforto. - F.}
\l{}
\l{\sou{entravar}. (\sou{trans.}) I. Impedar o jenar la marcho di animalo per}
\l{\cel{3}irga moyeno. - II. Impedar o jenar per ulo qua haltigas,}
\l{\cel{3}embarasas. - FS.}
\l{}
\l{\sou{entreo}. Disho qua venas pos la fisho-disho ed ante la rostajo.}
\l{\cel{3}- DFS.}
}
